% !TeX program = pdflatex
% papers/conversa_fundacao.tex — Corpus da escada (fundações), estilo conversa.
% Cresce como o casual: pares curtos; cita teoria.tex / torre_fundacao / tests.
\documentclass[11pt,a4paper]{article}
\usepackage[utf8]{inputenc}\usepackage[T1]{fontenc}\usepackage[brazil]{babel}
\usepackage[margin=2.4cm]{geometry}
\usepackage[hidelinks]{hyperref}
\newcommand{\code}[1]{\texttt{\small #1}}
% ─────────────────────────────────────────────────────────────────────────────
%  papers/gkcapa.tex — a capa do Reino, mínima, para os papers em `article`.
%
%  O estilo.tex completo é do `report` (títulos de capítulo, skak, …). Aqui fica
%  só o que a capa precisa, para o paper continuar a compilar sozinho com
%  pdflatex. O tradutor próprio (tests/tex) carrega o estilo.tex da raiz / o
%  symlink papers/estilo.tex e avalia o mesmo `\gkcapa`.
% ─────────────────────────────────────────────────────────────────────────────
\definecolor{ouro}{HTML}{B8912F}
\definecolor{ouroclaro}{HTML}{F3E9CE}
\definecolor{tinta}{HTML}{1E1B16}
\definecolor{regua}{HTML}{6E6858}
\providecommand{\gknota}{\fontsize{7.62}{11.05}\selectfont}
\providecommand{\gktexto}{\fontsize{10.50}{15.22}\selectfont}
\providecommand{\gksub}{\fontsize{12.33}{17.87}\selectfont}
\providecommand{\gksec}{\fontsize{14.47}{20.98}\selectfont}
\providecommand{\gkcap}{\fontsize{16.99}{24.63}\selectfont}
\providecommand{\gktit}{\fontsize{23.42}{33.95}\selectfont}
\providecommand{\Prj}{\mathbb{P}}
\providecommand{\Trf}{\mathcal{T}}
\providecommand{\gkcapa}[3]{%
\title{\vspace{-0.6cm}
{\color{ouro}\rule{\textwidth}{1.4pt}}\\[6mm]
\textbf{\gktit\color{tinta} Reino Dourado}\\[3mm]{\gksec\itshape\color{regua} Golden Kingdom}\\[8mm]
{\gkcap\scshape\color{ouro} #1}\\[5mm]
{\gksub\color{regua} #2}\\[2mm]
{\gktexto\color{regua}\itshape #3}\\[7mm]
{\color{ouro}\rule{\textwidth}{0.6pt}}\\[9mm]{\gknota\color{regua}\begin{minipage}{\textwidth}\setlength{\parindent}{0pt}%
\textbf{Aviso de Isenção Algorítmica:} O universo, as leis e as entidades contidas nestas páginas ---
incluindo felinos matriciais, aves de rapina algébricas e amuletos de controle de fluxo --- são
construções de pura ficção formal. Esta obra não descreve a realidade física, não serve como
engenharia reversa do cosmos e não constitui doutrina religiosa. Qualquer semelhança com bugs reais do seu
sistema operacional é mera coincidência (ou falta de reza). Prossiga por sua própria conta e
risco.\end{minipage}}}
\author{}\date{}}

\begin{document}
\gkcapa{Conversa Fundação}
       {Corpos $\to$ Dinâmica $\to$ Topologia $\to$ Hilbert/Clifford}
       {escada no corpus; Peano só depois}
\maketitle
\begin{abstract}
Pares de fala sobre as quatro fundações, no tom da assistente.
Fontes: \code{teoria.tex}, \code{papers/torre\_fundacao.tex},
\code{tests/topologia.c}, \code{tests/hurwitz.c}, \code{tests/hilbert\_bidual.c}.
Semente: \code{tools/fundacao.sh}.
\end{abstract}

\section{o que é um corpo}
Sítio com soma e produto onde os não-nulos invertem.
Ver \code{teoria.tex} e \code{papers/torre\_fundacao.tex}.

\section{o que é o corpo dual}
Par $(V,V^{*})$ com a identificação $J$ --- as leis vivem neste ambiente.
Detalhe em \code{teoria.tex}.

\section{o que é hurwitz}
Quatro álgebras de divisão normadas: $\mathbb{R},\mathbb{C},\mathbb{H},\mathbb{O}$.
Medido em \code{tests/hurwitz.c}; enunciado em \code{teoria.tex}.

\section{o que é dinamica de sistemas}
$X_{k+1}=F(X_k)$ --- órbitas, estados, volta.
\code{papers/torre\_fundacao.tex}; relógio em \code{teoria.tex}.

\section{o que é a topologia dos corpos}
Distância entre réguas: $d=\lvert\Delta_1-\Delta_2\rvert$, $\Delta=B^{2}-4C$.
\code{tests/topologia.c}.

\section{o que é a curva de hilbert aqui}
Deriva das Leis~1 e~2; bidual $\pi$ e $\nu$.
\code{tests/hilbert\_bidual.c}.

\section{o que é clifford}
Álgebra que orienta: $xy+yx=2\langle x,y\rangle$.
\code{papers/torre\_fundacao.tex}.

\section{mostra a fundação}
Escada: Corpos, Dinâmica, Topologia, Hilbert/Clifford --- depois Peano.
\code{papers/torre\_fundacao.tex}.

\section{o que é a lei 1}
$1^{\dagger}=-1$ --- dual/involução, período~2.
\code{teoria.tex} e \code{tests/hilbert\_bidual.c}.

\section{o que é a lei 2}
$T^{\dagger}=-T$ logo $T^{2}=-1$ --- rotor, período~4.
\code{teoria.tex} e \code{tests/hilbert\_bidual.c}.

\section{porque hurwitz para em 8}
Em 16 já há divisores de zero --- o cristal $N(xy)=N(x)N(y)$ parte.
\code{tests/hurwitz.c}.

\section{dois corpos sao isomorfos quando}
Quando $d=\lvert\Delta_1-\Delta_2\rvert=0$ --- mesma classe.
\code{tests/topologia.c}.

\section{o que é o bidual de hilbert}
$\pi$ estica e $\nu$ contrai; $\nu\circ\pi=\mathrm{id}$ e $\pi\circ\nu=\mathrm{id}$.
\code{tests/hilbert\_bidual.c}.

\section{o que é lyapunov}
Taxas de expansão/contração; dualizado $\lambda^{+}+\lambda^{-}=0$ na volta.
\code{papers/torre\_fundacao.tex}.

\section{o que é retracao}
$\pi\circ\iota=\mathrm{id}$ --- descer sem inventar.
\code{papers/torre\_fundacao.tex}.

\section{quando falamos do peano}
Depois das quatro fundações.
\code{papers/torre\_fundacao.tex} $\to$ \code{papers/corpo\_peano.tex}.

\section{o que é o teorema central}
Gentil$\leftrightarrow$Hurwitz pela medida --- \code{teoria.tex}; só depois da escada.

\end{document}
